\section{FluxOS application specification v9}
\label{sec:appspec-v9}

Every application that runs on \Flux is described by a small, owner-signed JSON document: a
container image, how much CPU, memory and disk it needs, how many copies to run, and for how long.
Because \FluxOS has no central scheduler and instead relies on every node independently deciding
what to run (\S9), that document is closer to a consensus artifact than to an ordinary
application-programming detail — every node on the network must parse and validate it the same
way, which is why changing its shape is slow and height-gated, exactly like a price table. This
section describes the specification version currently enforced on the network, and a substantial,
actively-developed redesign that is not enforced anywhere yet. This paper calls the specification
version \emph{v9} throughout, written in full here and as \specv{9} from this point on; it is
unrelated to \daemonversion, the version number of the consensus daemon \fluxd, which is
maintained independently and is never the referent of ``v9'' in this document
\srcintent{design intake}. The reason this section exists is structural, not archival: three fields
\specv{9} adds — \texttt{paymentMemo}, \texttt{referralCode} and \texttt{machineType} — are the
substrate that \PNR (\S7), the referral program, and the VPS/datacenter tier (\S11) are each stated
to depend on, and as of this writing none of the three has any consumer code anywhere in \FluxOS.

\subsection{Version gating as deployed}

\FluxOS accepts an incoming specification $s$ at chain height \hgt only if its declared version does
not exceed the maximum version enforced at that height, its top-level and per-component fields are
drawn from that version's whitelist, and it otherwise satisfies that version's own validity rules:
$$
\mathrm{Accept}(s,\hgt) \;\equiv\; v(s) \le V(\hgt) \;\wedge\; \mathrm{Fields}(s) \subseteq
\mathcal{F}_{v(s)} \;\wedge\; \mathrm{Valid}_{v(s)}(s),
$$
with $V(\hgt)$ the highest specification version whose enforcement height has been reached. This is
the same height-scheduling mechanism the network uses for price-table changes, applied to the
specification schema itself: \texttt{verify\allowbreak{}App\allowbreak{}Specifications} runs type correctness, then
per-version restriction checks, then the enforcement-height gate, then object-key correctness
against a per-version field whitelist, in that fixed order, before anything is broadcast \shipped\ \srccode{flux/\allowbreak{}ZelBack/\allowbreak{}src/\allowbreak{}services/\allowbreak{}appRequirements/\allowbreak{}appValidator.js}{1243-1400}.

The enforcement-height table, as configured today, is:

\begin{table}[H]
\centering
\small
\begin{tabular}{@{}llp{6.5cm}@{}}
\toprule
Version & Enforcement height & Notes \\
\midrule
1 & 0 & no longer accepted for registration or update via the API \\
2 & 0 & \\
3 & 983{,}000 & \\
4 & 1{,}004{,}000 & composition (\texttt{compose} array) \\
5 & 1{,}142{,}000 & \texttt{contacts}, \texttt{geolocation} \\
6 & 1{,}300{,}000 & \texttt{expire}, price change \\
7 & 1{,}420{,}000 (mainnet) / 1{,}390{,}000 (dev) & \texttt{nodes} targeting, \texttt{secrets}, \texttt{staticip} \\
8 & \textbf{1{,}932{,}380} (mainnet) / \textbf{1{,}921{,}500} (dev) & enterprise / ArcaneOS apps \\
\bottomrule
\end{tabular}
\caption{Specification-version enforcement heights as shipped in \FluxOS. There is no row for
version 9: no enforcement height for it exists in the shipped configuration or the live pricing
endpoint (below).}
\end{table}

\shipped\ \srccode{flux/\allowbreak{}ZelBack/\allowbreak{}config/\allowbreak{}default.js}{231-240}, cross-checked against the live public
pricing endpoint, which returns the identical mainnet heights through version 8 and carries no key
for version 9 at all \srcmeasured{api.runonflux.io/\allowbreak{}apps/\allowbreak{}deploymentinformation}{2026-09-03}.
\specv{8}'s enforcement height is \textbf{1{,}932{,}380} on mainnet and \textbf{1{,}921{,}500} in the
development network \shipped\ \srccode{flux/\allowbreak{}ZelBack/\allowbreak{}config/\allowbreak{}default.js}{231-240}. Above that height, a
specification whose top-level keys are anything other than \texttt{version,\allowbreak{} name,\allowbreak{} description,\allowbreak{}
owner,\allowbreak{} compose,\allowbreak{} instances,\allowbreak{} contacts,\allowbreak{} geolocation,\allowbreak{} expire,\allowbreak{} nodes,\allowbreak{} staticip,\allowbreak{} datacenter,\allowbreak{} enterprise}
is rejected outright \shipped\ \srccode{flux/\allowbreak{}ZelBack/\allowbreak{}src/\allowbreak{}services/\allowbreak{}appRequirements/\allowbreak{}appValidator.js}{1063-1067}.

\inprogress\ The organisation's own working checkout additionally carries a single, unpublished
commit that begins to vendor a copy of the RFC-0 schema (\S\ref{sec:two-artifacts}) behind a feature flag
rather than a height: registering or updating a \texttt{version: 9} specification is rejected
before schema validation ever runs unless \texttt{config.\allowbreak{}fluxapps.\allowbreak{}acceptV9} is \texttt{true}, and
that flag defaults to \texttt{false} with no caller anywhere in the codebase that sets it
\srcbranch{RunOnFlux/\allowbreak{}flux}{feat/appspec-v9}{ZelBack/src/services/appDatabase/registryManager.js:1766-1769}.\footnote{%
This sits on a working checkout of the organisation repository that had not been pushed to
\texttt{RunOnFlux/flux}'s public remote at the time of writing, and it is a much smaller and
separate effort from the contributor fork discussed in \S\ref{sec:appspec-v9-scale}. It is cited
here as an unpublished working branch, not as a fetchable public artifact.} No node on the network
today enforces any rule for \texttt{version: 9} beyond ``reject it.''

\subsection{The migration baseline, measured}
\label{sec:migration-baseline}

The concrete cost of this gating scheme is visible in what is actually deployed. Of 1{,}195
applications registered on the network on 2026-09-03
\srcmeasured{api.runonflux.io/\allowbreak{}apps/\allowbreak{}globalappsspecifications}{2026-09-03}:

\begin{table}[H]
\centering
\begin{tabular}{@{}lrr@{}}
\toprule
Version & Applications & Share \\
\midrule
v8 & 873 & 73.1\% \\
v7 & 209 & 17.5\% \\
v6 & 70 & 5.9\% \\
v3 & 27 & 2.3\% \\
v4 & 11 & 0.9\% \\
v5 & 3 & 0.3\% \\
v2 & 2 & 0.2\% \\
\midrule
Total & 1{,}195 & 100\% \\
\bottomrule
\end{tabular}
\caption{Deployed-application count by specification version, ordered by share.}
\end{table}

\shipped\ \textbf{322 applications (26.9\%) are on a version older than \specv{8}.} Seven
specification versions — \specv{2} through \specv{8} — are live on the network simultaneously
\srcmeasured{api.runonflux.io/\allowbreak{}apps/\allowbreak{}globalappsspecifications}{2026-09-03}. The acceptance predicate
above has a ceiling ($v(s) \le V(\hgt)$) but no floor: nothing in it forces an already-running
application to move forward when a newer version becomes enforceable, and the measured coexistence
of seven versions is the direct, observable consequence. The network has not deprecated any
specification version to date; every version ever shipped is still validated, by every node, for as
long as an application on that version continues to exist. That is the real cost of
height-scheduled versioning with no deprecation mechanism, and it is the baseline against which
\specv{9}'s eventual migration has to be measured (\S\ref{sec:migration-mechanism}).

\subsection{What v9 changes}

\inprogress\ \specv{9} is defined, at present, only as a draft JSON Schema (RFC-0) in a standalone
repository \srccode{flux-spec/\allowbreak{}schema/\allowbreak{}v9.json}{1-350} — see \S\ref{sec:two-artifacts} for why that
draft and what \FluxOS actually consumes are not the same object. Table~\ref{tab:v9-changes}
summarises the delta against \specv{8}: four field renames, one removal, and five additions, each
against what depends on it and how far its implementation actually reaches.

\begin{table}[H]
\centering
\small
\begin{tabular}{@{}llp{3.6cm}p{4.4cm}@{}}
\toprule
Change & \specv{9} name (\specv{8} name) & Depends on & Implementation status \\
\midrule
Rename & \texttt{components} (\texttt{compose}) & --- & schema-level only; breaks every existing client, hence bundled with the other renames rather than dripped \\
Rename & \texttt{image} (\texttt{repotag}) & --- & schema-level only \\
Rename & \texttt{memory} (\texttt{ram}) & --- & schema-level only \\
Rename & \texttt{rootFsGb} (\texttt{hdd}) & --- & schema-level only \\
Removal & --- (\texttt{enterprise}; its envelope replacement named but not specified) & private/encrypted specifications & regression; see \S\ref{sec:envelope-regression} \\
Addition & \texttt{loadBalancing} & routing, health checks, backend TLS, connection draining & partially implemented: draining and per-app-CA TLS wired; config generation stays in \FDM (\S12) \\
Addition & \texttt{duration} & subscriptions and auto-renewal (\S16) & first-class in the schema; the fork's pricing engine reads the same concept off \texttt{spec.ttl}, and whether \texttt{duration} maps to it is unresolved (\S\ref{sec:app-interfaces}) \\
Addition & \texttt{paymentMemo} & \PNR (\S7) & \textbf{zero consumers in \FluxOS} \\
Addition & \texttt{referralCode} & referrals, preview environments & \textbf{zero consumers} \\
Addition & \texttt{machineType} & VPS/datacenter tier (\S11) & \textbf{zero consumers} \\
\bottomrule
\end{tabular}
\caption{\specv{8}$\to$\specv{9} field changes, their dependents, and implementation status.}
\label{tab:v9-changes}
\end{table}

Field shapes and required-field membership per \srccode{flux-spec/\allowbreak{}schema/\allowbreak{}v9.json}{40-46} (\texttt{components}),
\srccode{flux-spec/\allowbreak{}schema/\allowbreak{}v9.json}{47-61} (\texttt{loadBalancing}), \srccode{flux-spec/\allowbreak{}schema/\allowbreak{}v9.json}{62-82}
(\texttt{duration}, \texttt{paymentMemo}, \texttt{referralCode}, \texttt{machineType}), and
\srccode{flux-spec/\allowbreak{}schema/\allowbreak{}v9.json}{103-122} (\texttt{image}, \texttt{memory}, \texttt{rootFsGb}); the
removal of \texttt{enterprise}, with its envelope replacement deferred to a later revision, is stated in the draft's own README
\srcdoc{flux-spec/README.md:58-67,95-96}; the \texttt{paymentMemo}\slash\texttt{referralCode} dependency
mapping is drawn from the fork's own code comments (``basis for paying operators under PNR v2'',
``rewards paid as hosting credits'') \srcbranch{RunOnFlux/\allowbreak{}flux}{feat/appspec-v9}{ZelBack/src/services/appRequirements/appSpecV9.js}.

\subsection{The finding this section exists to deliver}
\label{sec:finding}

\inprogress\ Stated plainly and without spin: \texttt{paymentMemo}, \texttt{referralCode} and
\texttt{machineType} have \textbf{zero consumer code anywhere in \FluxOS} — no billing path, no
referral-reward path, no
capacity-tier scheduling path reads any of the three, confirmed by grepping every JavaScript file
under \texttt{ZelBack/} and \texttt{tests/} for all three field names
\srcbranch{MorningLightMountain713/\allowbreak{}flux}{feat/fluxos-v9}{(grep -rn "paymentMemo|referralCode|machineType" --include=*.js ZelBack tests)}.
Those are exactly the fields \PNR (\S7), the referral program, and the VPS tier (\S11) are stated to
depend on. The orchestration substrate of \specv{9} — placement, mesh identity, routing plumbing,
persistent storage, pricing — is built and tested (\S\ref{sec:appspec-v9-scale}); the economic
fields the roadmap's dependency chain actually rests on are declared in a schema and nothing more.
The dependency chain the roadmap draws between \specv{9} and \PNR, referrals and the VPS tier is
real at the level of the schema; it currently terminates in fields no code reads. \S\ref{sec:feasibility} returns to
this as one instance of a broader pattern in how far the remaining work actually reaches.

\subsection{The scale of the work, quantified}
\label{sec:appspec-v9-scale}

The overwhelming majority of \specv{9} engineering effort does not live in the organisation
repository. It lives on an individual contributor's fork\footnote{\texttt{MorningLightMountain713/\allowbreak{}flux},
branch \texttt{feat/fluxos-v9}, last pushed 2026-09-02 — an external fork of \texttt{RunOnFlux/flux},
not a branch of the organisation repository itself. Citations against it use
\textsf{[repo@branch:file:line]} notation for that reason, and none of them can be fetched from the
organisation's own remote.} that diffs 932 files against \texttt{upstream/\allowbreak{}master}, with
$+180{,}317/-49{,}561$ lines
\srcbranch{MorningLightMountain713/\allowbreak{}flux}{feat/fluxos-v9}{(git diff --stat upstream/master origin/feat/fluxos-v9)},
decomposing the old monolithic application-service layout into 23 focused service
packages, ten of them absent from \texttt{upstream/\allowbreak{}master} (\texttt{appMesh}, \texttt{appPlacement}, \texttt{pricing},
\texttt{quorumGrant}, and others)
\srcbranch{MorningLightMountain713/\allowbreak{}flux}{feat/fluxos-v9}{(find ZelBack/src -maxdepth 3 -type d)}.
This is \inprogress\ work, and it is substantial: 289 unit-test files containing 7{,}763 test cases
run on every push
\srcbranch{MorningLightMountain713/\allowbreak{}flux}{feat/fluxos-v9}{(find tests/unit -name *.test.js | wc -l = 289; grep count of it/test call sites in tests/unit = 7763)}
\srcbranch{MorningLightMountain713/\allowbreak{}flux}{feat/fluxos-v9}{.github/workflows/nodejs.yml:1-3,90-93},
alongside a from-scratch, 121-scenario integration harness that provisions up to 16 simulated nodes
with per-daemon stubs
\srcbranch{MorningLightMountain713/\allowbreak{}flux}{feat/fluxos-v9}{(ls test-infra/runner/tests); (find test-infra/config -maxdepth 1 -type d)}
— but that harness runs only on manual dispatch, on a self-hosted runner, and is not part of the
automatic merge gate
\srcbranch{MorningLightMountain713/\allowbreak{}flux}{feat/fluxos-v9}{.github/workflows/integration-harness.yml:1-38}.

This is an ordinary open-source pattern, not a criticism, but it is a fact about where development
authority currently sits for the largest in-flight change to how applications are described on
Flux: on a contributor's fork rather than in the organisation's own repository. \S23 returns to
this as a governance observation.

\subsection{The routing layer, precisely}

\srcroadmap{flux-roadmap-public.md:52-53} The public roadmap states that \specv{9}'s routing layer
has been ``rewritten.'' Read precisely against the fork, that claim is partially, not wholly,
supported. \inprogress\ Two mechanisms declared through the \texttt{loadBalancing} block are
genuinely wired: connection draining, served over a UNIX-socket
JSON-RPC endpoint that \texttt{flux-shutdownd} calls with \texttt{drain\_app}\slash\texttt{stop\_app} to
mark an application draining ahead of node shutdown
\srcbranch{MorningLightMountain713/\allowbreak{}flux}{feat/fluxos-v9}{ZelBack/src/services/appMessaging/drainServer.js:12-24,85-136,176-201};
and per-application-CA backend TLS, under which a node generates a local Ed25519 keypair, obtains a
30-day host certificate signed by the application's own certificate authority through \fluxbench,
and renews it automatically before expiry
\srcbranch{MorningLightMountain713/\allowbreak{}flux}{feat/fluxos-v9}{ZelBack/src/services/appLifecycle/backendTlsService.js:1-18,73-116}.
What is \emph{not} true is that \FluxOS now owns HAProxy configuration generation: grepping every
service file for the routing-layer vocabulary (\texttt{loadBalancing}, \texttt{healthCheck},
\texttt{stickySession}, \texttt{connection\allowbreak{}Drain}, \texttt{customDomain}) turns up three files, none
of them a route-table or config-generation engine
\srcbranch{MorningLightMountain713/\allowbreak{}flux}{feat/fluxos-v9}{(grep -rln "loadBalancing|healthCheck|stickySession|connectionDrain|customDomain" ZelBack/src/services)}.
That responsibility stays where it already sits, in \FDM (\S12). And the actual signal a routing
layer would consume when an application starts draining is not a new push protocol; marking a
component draining triggers an immediate presence rebroadcast, so the mechanism by which \FDM
would pull a backend out of rotation is the same gossip/presence-announcement path \FluxOS already
used before \specv{9}
\srcbranch{MorningLightMountain713/\allowbreak{}flux}{feat/fluxos-v9}{ZelBack/src/services/appMessaging/drainServer.js:79-91}.
So: the spec-side declaration and two real backend mechanisms are new and wired; the config-generation
engine the roadmap's phrase might suggest is not, and never was scoped to live here.

\subsection{Two artifacts, not one}
\label{sec:two-artifacts}

``\specv{9}'' currently names two different objects, and this section describes them separately
rather than presenting one as a stand-in for the other.

The first is \texttt{Run\allowbreak{}On\allowbreak{}Flux/\allowbreak{}flux-\allowbreak{}spec}, a standalone repository holding the RFC-0 draft schema:
two commits, both dated 2026-09-01, self-described as ``Status: RFC-0 — draft for comment
(2026-09-01). Not live on the network'' \srcdoc{flux-spec/README.md:3}, unpublished
(\texttt{"private":\allowbreak{} true}, no \texttt{bin}\slash\texttt{main}\slash\texttt{exports}), with no CI and no tags.
The second is the vendored copy the contributor fork actually consumes at runtime, which is pulled
not from this repository but from a different, private, monorepo-structured \texttt{flux-spec}
pinned at a specific commit hash
\srcbranch{MorningLightMountain713/\allowbreak{}flux}{feat/fluxos-v9}{test-infra/flux-spec/pin:1}; the two are
structurally different (the private monorepo has \texttt{packages/spec}, \texttt{spec-backend},
\texttt{spec-policy} subpackages that do not exist in the public draft) and, on the evidence
gathered, disagree on at least one field name — the draft's \texttt{duration}
\srccode{flux-spec/\allowbreak{}schema/\allowbreak{}v9.json}{15,62} versus the fork's internal \texttt{spec.ttl}
\srcbranch{MorningLightMountain713/\allowbreak{}flux}{feat/fluxos-v9}{ZelBack/src/services/pricing/v9PricingRegime.js:95,108,122,128,183,228,279,308}
— which this paper flags rather than resolves, since the vendored validator classes themselves
are not available to inspect.

\inprogress\ The public draft carries two specific defects worth naming, because a specification
whose entire purpose is that every node agrees on it should not have them. First, it has \textbf{no signature
field at all}: nothing in \texttt{schema/v9.json} names a field a submitter signs, and the only
``signature''/``signed'' references anywhere in the repository describe an unrelated \fluxbench
attestation workstream \srccode{flux-spec/\allowbreak{}schema/\allowbreak{}v9.json}{1-350}. A ``recommended'' canonicalize-then-
double-SHA256 fingerprint recipe exists for comparison purposes, but the draft offers it only as a
fingerprint and nowhere defines what gets signed for submission
authenticity \srcdoc{flux-spec/validation/canonicalization.md:53-63}. Second, the draft's own
canonicalization rules — array-sort ordering for fields such as \texttt{customDomains} and
\texttt{certificate\allowbreak{}Authorities}, and explicit materialization of defaulted fields — are stated as
\texttt{MUST} requirements in prose, with \textbf{no schema or code enforcement} anywhere in the
repository \srcdoc{flux-spec/validation/canonicalization.md:31-45}. JSON Schema has no array-sort
keyword, so nothing here can express these rules as validation; a document can be schema-valid and
canonically invalid at the same time, and nothing in the draft detects that. For a document meant to
be compared byte-for-byte across roughly 6{,}500 independent validators, that is the defect to name.

\subsection{The encryption-envelope regression}
\label{sec:envelope-regression}

\planned\ The public roadmap states that \specv{9} will deliver ``a cleaner encryption envelope for
private specs'' \srcroadmap{flux-roadmap-public.md:61}. \inprogress\ The RFC-0 draft does the
opposite of deliver it: it removes \specv{8}'s \texttt{enterprise} field outright, names the
encryption envelope as its replacement \textbf{without specifying it}, and
lists that envelope for private deployments as an explicitly open question deferred to a
future schema revision \srcdoc{flux-spec/README.md:58-67,95-96}. Concretely, \specv{9} as drafted
has no way to express a private or enterprise deployment at all. This is not a marginal population:
223 of the 1{,}195 deployed applications measured in \S\ref{sec:migration-baseline} (18.7\%) use
\texttt{enterprise} today \srcmeasured{api.runonflux.io/\allowbreak{}apps/\allowbreak{}globalappsspecifications}{2026-09-03}.
If the published draft is what ships, the roadmap's promise is currently unmet in it, and this gap
has to be closed before \texttt{enterprise} can be removed on any node (\S\ref{sec:migration-mechanism}).

\subsection{The migration mechanism: forced, not opt-in}
\label{sec:migration-mechanism}

\inprogress\ No mechanism for migrating deployed applications from \specv{8} (or older) to \specv{9}
exists in code today: the fork's own \texttt{migrations/} service registers exactly two node-runtime-state
migrations (node identity, container labels), neither of which touches application specifications
\srcbranch{MorningLightMountain713/\allowbreak{}flux}{feat/fluxos-v9}{ZelBack/src/services/migrations/index.js:39-51}.

\planned\ What is decided supersedes an earlier draft of this subsection, which proposed a
non-destructive version transition of this paper's own design: a version floor $v_{\min}$ with a
deprecation schedule, below which existing deployments would be refused renewal rather than
evicted. The team's design is a different one: migration to \specv{9} will be
\textbf{forced, not opt-in}. Every application will be migrated by an application-update message
carrying a special signature that compels the \specv{9} version; updates and upgrades on older
specification versions will then be disabled, and once the entire fleet is on \specv{9} the legacy
code paths will be removed \srcintent{design intake}. This paper states
the change rather than silently substituting it: the forced design is operationally cleaner than the
deprecation window it replaces, because it would remove the 322-application, 26.9\% legacy-version
tail (\S\ref{sec:migration-baseline}) in a single step, at the moment the forcing message is
broadcast, rather than letting it trail off across a schedule at each owner's own pace.

The invariant this section stated for the superseded design was that no deployed application is
evicted by a version change — it lapses under the rules it was created under. Forced migration is
intended to preserve non-eviction by a different route, and the difference is worth stating
precisely, as a design target rather than a shipped guarantee:

\begin{property}
No application is evicted by the migration to \specv{9}: instead of lapsing under the rules it was
created under, every application is rewritten forward onto the new version.
\end{property}

\planned\ What this would trade against the superseded invariant is consent: nothing in the design
gives an owner a way to remain on their signed specification, or to decline the rewrite
\srcintent{design intake}. Non-eviction would be preserved; opt-out
would not exist.

A key able to rewrite any tenant's application specification without that tenant's own signature
would be a privileged capability, and \srcinferred\ this paper places it, on that basis, in the same
class it already inventories elsewhere: the unsigned network policy (\S12), the OS release channel
(\S14), and the 2-of-4 emergency block keys (\S4). Naming it is the point; a capability of this
shape left unnamed is a finding, not a footnote. \planned\ Per the team's decision, three properties
are intended to bound it: it would be \textbf{single-use}, scoped to the migration rather than a
standing signing service; \textbf{migration-scoped}, existing only to perform this one version
transition and nothing else; and \textbf{hardcoded}, expected to be baked-in values rather than a
live signer that could be invoked again later
\srcintent{design intake}. Those three properties are what would make the
capability acceptable rather than alarming.

\planned\ The scope and the enforcement class are now decided rather than open. The scope is the
\textbf{version field and the required-field additions only} --- never the container image and never
the application's functionality --- and the stated purpose is strictly to bring the network into
alignment on one schema version, nothing more
\srcintent{design intake}. Enforcement is by \textbf{\FluxOS}, not by
\fluxd: every node validates the forced-update message against that scope, which is the same
enforcement class as specification validation generally --- the \texttt{verify\allowbreak{}App\allowbreak{}Specifications}
check described above --- checked by every node, but as a \FluxOS{} policy rule rather than a
consensus rule. That distinction bounds the capability's blast radius precisely: a future \FluxOS{}
release could in principle change what the forcing message is permitted to touch, in exactly the
sense that \S23's deployed policy channel is \FluxOS-enforced and not consensus state, and the two
capabilities belong to the same class for that reason. It also means the scope is a stated design
constraint rather than a proof: nothing in \fluxd{} itself rejects a forced-update message that
exceeded it, so the guarantee rests on the migration tooling matching its own specification rather
than on a rule \fluxd{} would enforce independently. Explicit retirement of the key once the
migration completes remains a step the team has not stated as a commitment, and leaving hardcoded
values live in a binary indefinitely after that point would be a departure from the single-use
property above.

\planned\ Under this design the four field renames (\texttt{compose}$\to$\texttt{components},
\texttt{repotag}$\to$\texttt{image}, \texttt{ram}$\to$\texttt{memory},
\texttt{hdd}$\to$\texttt{rootFsGb}) would need no separate compatibility shim: the same
forced-update message that compels the version bump would rewrite the renamed fields as part of the
same rewrite, so the translation question this section previously posed as open is resolved by the
mechanism itself rather than by added code.

\subsection{The v9 schema, for reference}

\inprogress\ Listing~\ref{lst:v9-schema} reconstructs the root object of the draft schema from the
field table gathered during evidence collection; it is not a verbatim excerpt of the source file;
the citations that follow it, not the listing itself, are the evidence that these fields exist.

\begin{lstlisting}[caption={Root object of the \specv{9} RFC-0 draft schema (reconstructed). Look at \texttt{additional\allowbreak{}Properties:\allowbreak{}false} together with the five new required fields: this combination is what makes an unrecognised key --- including v8's removed \texttt{enterprise} --- a hard validation failure rather than a silent drop, and what makes \texttt{paymentMemo}\slash\texttt{referralCode}\slash\texttt{machineType} required inputs to a document that \S\ref{sec:finding} shows no \FluxOS code yet reads.}, label={lst:v9-schema}]
{
  "type": "object",
  "additionalProperties": false,
  "required": ["version","name","description","owner",
               "components","loadBalancing","duration",
               "paymentMemo","referralCode","machineType"],
  "properties": {
    "version":      { "const": 9 },
    "duration":     { "type": "integer", "minimum": 1 },
    "paymentMemo":  { "type": "string", "minLength": 1, "maxLength": 512 },
    "referralCode": { "type": "string", "minLength": 1, "maxLength": 64 },
    "machineType":  { "enum": ["standard", "gpu", "vps"] }
  }
}
\end{lstlisting}

Root object, \texttt{additional\allowbreak{}Properties:\allowbreak{}false}, and the ten required top-level fields:
\srccode{flux-spec/\allowbreak{}schema/\allowbreak{}v9.json}{7-19}. \texttt{version}: \srccode{flux-spec/\allowbreak{}schema/\allowbreak{}v9.json}{21-24}.
\texttt{duration}: \srccode{flux-spec/\allowbreak{}schema/\allowbreak{}v9.json}{62-66}. \texttt{paymentMemo}:
\srccode{flux-spec/\allowbreak{}schema/\allowbreak{}v9.json}{67-72}. \texttt{referralCode}: \srccode{flux-spec/\allowbreak{}schema/\allowbreak{}v9.json}{73-78}.
\texttt{machineType}: \srccode{flux-spec/\allowbreak{}schema/\allowbreak{}v9.json}{79-82}.

\subsection{Open parameters}

\begin{remark}[\specv{9} timing: the quarter settled in round 11, the height in round 17]
\label{rem:v9-timing}
\specv{9} is in development and is to be \textbf{activated in Q4~2026, with the migration of
existing applications in the same quarter} \srcintent{design intake} (\planned).
That places $\hstar$ before $\hpa = \num{3787502}$ (2027-07-01), so the specification change
lands ahead of \PNR{} activation and the parallel-asset cut rather than with them --- which is the
right order, because \S\ref{sec:feasibility} places \PNR{}, credits, referrals, the VPS tier and the
provisioning API all downstream of the \specv{9} economic fields.

Round 11 fixed no height, and the constraint on it is unchanged: $\hstar$ must \emph{follow} the
implementation of the \texttt{paymentMemo}\slash\texttt{referralCode}\slash\texttt{machineType} trio
(\S\ref{sec:finding}), never precede it, or the height would activate a schema whose economic fields
are declared but inert. At the \SI{30}{\second} target a Q4~2026 activation corresponds to roughly
$\num{3000000}$--$\num{3260000}$, and the paper states that range as an implication of the stated
quarter rather than as a number the team has given.
\settlednote{$\hstar = \num{3050000}$ \srcintent{design intake},
approximately 2026-10-20 at the \SI{30}{\second} target --- inside the settled Q4~2026 window, about
\num{45} days beyond the 2026-09-05 tip of \num{2921193}, and about \num{256} days before
$\hpa = \num{3787502}$. That ordering is what \S\ref{sec:feasibility} requires: the \specv{9}
economic fields land, and are given roughly eight months to prove themselves, before \PNR{}, credits,
referrals, the VPS tier and the provisioning API activate on top of them. The constraint of
\Cref{rem:v9-timing} still binds and is now a scheduling obligation rather than an open question:
the \texttt{paymentMemo}\slash\texttt{referralCode}\slash\texttt{machineType} implementation must be
complete before this height, or the height activates a schema whose economic fields are inert.}
\end{remark}

\settlednote{\texttt{quorum\allowbreak{}Grant\allowbreak{}Activation\allowbreak{}Height} activates \emph{after} \specv{9}, once
\specv{9} is stable, at a height not fixed now
\srcintent{design intake}. This is a deliberate decoupling rather than an
unfinished decision: the mastership/active-standby quorum plane and the \specv{9} schema change have
unrelated failure modes, so binding them to one height would mean a fault in either forces a rollback
of both. The cost is a second upgrade event and a second announcement. The height is left unset
because the trigger is an observed condition --- \specv{9} running stably --- and a height chosen
today could only be a guess at when that condition will hold. Until it activates, the field remains
what \S\ref{sec:one-platform} measures it as: read by production code, set only in integration-test
fixtures, and inert on mainnet.}

%% Drain this section's float queue before the next begins.
\FloatBarrier
